\section{Experiments}

The theory is checked on a bench, a finite, discrete, deterministic, reproducible simulator that turns each question into one experiment returning a graded verdict. The catalog holds over four hundred such experiments across eighteen families, and the strength of the program lives in the method rather than in any one result.

\textbf{The depth rubric.} Each verdict carries a grade. L0 is circular, the answer put in by hand, kept only as a consistency note. L1 reproduces a known mathematical fact on the substrate. L2 reproduces a known physical construction on the substrate. L3 is an emergent and novel result, with a control that could have failed. Only L3 is both new and guarded, and only a handful reach it cleanly, chief among them the emergent renormalized rule that beats its naive control and the integrated-information measure that ranks a real self far above a random bag.

\textbf{The headline results.} Table~\ref{tab:reproduced} gathers what the substrate reproduces, each row the outcome of one or more experiments in the catalog, every number a measurement of the integer dynamics rather than a constant placed into the law by hand.

\begin{table}[ht]
\centering
\small
\begin{tabular}{@{}ll@{}}
\toprule
what the substrate reproduces & the value it lands on \\
\midrule
a clean chiral fermion & spin one-half, no lattice doubling \\
the gyromagnetic ratio $g$ & $2$ \\
generations of matter & three \\
the unifying gauge group & $\mathrm{SO}(10)$ down to the Standard Model \\
the weak mixing angle $\sin^2\theta_W$ & $3/8$ at unification, running toward $0.231$ \\
spatial dimensions selected & three \\
the spectral dimension of space & $\approx 2.97$ \\
the light cone and Lorentz symmetry & emergent, the flat control failing \\
gravity in the cusp & Newton's law and the Einstein equations \\
light bending by a mass & twice the Newtonian value \\
black-hole entropy & the area, not the volume \\
cosmic expansion & accelerating, $H = \tfrac{1}{3}\ln(\text{growth})$ \\
the quantum correlation (CHSH) & the Tsirelson bound $2\sqrt 2$ \\
the amplitude-to-probability law & the Born square, exponent two \\
\bottomrule
\end{tabular}
\caption{A summary of what the discrete substrate reproduces, each entry a measurement of the integer dynamics rather than a number placed into the law by hand.}
\label{tab:reproduced}
\end{table}

\textbf{The families.} The catalog is sorted into families, each gathering the questions of one domain. A first group tests the substrate itself, the \textit{foundations} that fix the dock and the tones, the \textit{geometry} of the honeycomb and its growth, the \textit{substrate survey} that weighs the committed tessellation against its rivals, and the \textit{spin} structure that singles it out. A second group tests the machinery of the law, the \textit{general} behavior of the knit, the \textit{data structures} and \textit{addressing} it supports, the \textit{associative} recall that rides on them, the \textit{computation} it can carry, and the \textit{renormalization} that climbs from fine grain to coarse. The physics families take on the \textit{gauge} fields and forces, \textit{gravity}, \textit{relativity} and the light cone, the \textit{quantum} correlations, \textit{holography}, and \textit{cosmology}. And the \textit{selves} family tests whether bound, self-correcting subjects form at all and how strongly they integrate. Table~\ref{tab:families} gives one representative question and its grade for each.

\begin{table}[ht]
\centering
\small
\begin{tabular}{@{}lll@{}}
\toprule
family & a representative question & grade \\
\midrule
foundations & is the dock forced by four independent constructions? & L1 \\
geometry & does the crystal grow at the hyperbolic rate? & L1 \\
addressing & does every dock carry a solved coordinate? & L1 \\
spin & does the directional state carry a true spinor? & L2 \\
quantum & does the CHSH correlation reach $2\sqrt 2$? & L2 \\
holography & does boundary entropy count bulk area? & L2 \\
computation & can the substrate emulate a universal machine? & L2 \\
associative & does a fragment recall its stored whole? & L2 \\
gauge & does the weak mixing angle run to its measured value? & L3 \\
relativity & does a finite light cone and Lorentz symmetry emerge? & L3 \\
gravity & does a bound self feel an effective Newtonian pull? & L3 \\
cosmology & does the universe accelerate and select three dimensions? & L3 \\
renormalization & does the coarse rule beat its naive control? & L3 \\
selves & does an integrated self rank far above a random bag? & L3 \\
\bottomrule
\end{tabular}
\caption{One representative question per family, with the depth grade its cleanest result reaches.}
\label{tab:families}
\end{table}

\textbf{The discipline.} A result counts only when its control could have come out the other way. The model is deterministic, never reaching for a random seed, and where a trend is wanted it is found by growing the run rather than averaging over draws. Every claim must trace back to the eight elements with no ninth thing smuggled in, so a phenomenon that does not emerge is reported as a negative rather than forced by a quiet addition. Measuring a hyperbolic bulk has its own hazards, since a finite patch is almost all boundary, and the bench meets them with Bethe recursion, closed manifolds, band theory, and kernel-polynomial methods that read bulk-local quantities. Results are sorted into three layers by what they depend on, bulk findings that hold on any hyperbolic substrate, boundary findings that need the curvature, and flat-layer findings that depend on the clean cubic cusp.

\textbf{Controls and negative results.} The program publishes its failures. The empty vacuum exerts no force at a distance. A naive radial coarse-graining produces no tower of scales. No single committed rule both confines a body and radiates from it. A self cannot fully store itself. A like-charged pair of selves repels rather than binds. Each is kept in view as part of the case rather than hidden from it.

\textbf{The source of truth.} The codebase is the authority for everything stated here, and a complete list of every experiment with its current status lives in the codebase data alongside the runnable bench, so each verdict in this appendix can be traced to a specific test and rerun directly. \cite{pollard2026vibetest}
